\section{Affine square groups}\label{affine-square-groups}

Let \[
 a(x,y)=(-x,y),\quad b(x,y)=(y,x),\quad c(x,y)=(x,1-y).
\] \textbf{Proposition UA.Affine (the complete square presentation).}
These generate \(W=\mathbb Z^2\rtimes D_8\), of type \(\{4,4\}\), with
\(ac=ca\). Its presentation is \[
 W=\langle a,b,c\mid a^2=b^2=c^2=(ab)^4=(bc)^4=(ac)^2=1\rangle. \tag{S1}
\] \emph{Proof.} Put \(d=bab\), \(y_0=cd\), and \(x_0=by_0b\). The
relations give \([a,d]=1\), \([a,y_0]=1\), \(dy_0d=y_0^{-1}\),
\(ax_0a=x_0^{-1}\), and conjugation by \(b\) swaps \(x_0,y_0\).
Substitution of \(bc=x_0ab\) makes \((bc)^4=1\) equivalent to
\([x_0,y_0]=1\). Thus every word collects to \(x_0^u y_0^v M\),
\(M\in\langle a,b\rangle=D_8\). In the affine model \(x_0,y_0\) are the
unit translations, so these forms are distinct. This proves (S1).
\(\square\)

The normal form also fixes our affine multiplication convention: the
rightmost transformation acts first. The symbol \(D_8\) denotes the
group of signed permutation matrices in two coordinates, of order eight.
In all dihedral notation below the subscript denotes the order of the
group.

\textbf{Proposition UA.Square (finite square seeds).} For \(n\ge5\),
define the invariant lattice \[
 \Lambda_n=\begin{cases}
  2^e\mathbb Z^2,&n=2e+3,\\
  2^e\langle(1,1),(1,-1)\rangle,&n=2e+4,
 \end{cases}\qquad e\ge1,
\] and put \(H_n=W/\Lambda_n\). Its order is \(2^n\); the full linear
\(D_8\) coordinate is retained, even if its action on the vertex
quotient is not faithful. Its pair orders are \(4,4,2\). It has all
Boolean deletions and IP. Its third lower-central subgroup is
\(2\mathbb Z^2/\Lambda_n\), of order \(2^{n-5}\).

\emph{Proof.} The lattice indices are respectively \(2^{2e}\) and
\(2^{2e+1}\), so the normal form gives the stated orders. The three
linear pair products have orders four, four and two; hence the marked
pair orders cannot fall in either lattice quotient.

Write \(T=abcb\), \(U=cbab\), translations by \((-1,0)\), \((0,1)\). The
extra relators for \(H_n\) can be taken as \(T^{2^e},U^{2^e}\) in the
square case, and \(T^{2^e}U^{2^e},T^{2^e}U^{-2^e}\) in the diagonal
case. After deleting any marker, both \(T,U\) have order at most two:
deleting \(a\) leaves conjugates of \(c\), deleting \(c\) leaves
conjugates of \(a\), and deleting \(b\) leaves \(ac,ca\). Therefore
every lattice relator survives every deletion. The even Coxeter relators
also survive. Thus \(H_n\) has Boolean deletions and IP by UA.D.

Moreover, \[
 \gamma_3(W)=2\mathbb Z^2,
 \qquad \gamma_3(H_n)=2\mathbb Z^2/\Lambda_n,
 \qquad |\gamma_3(H_n)|=2^{n-5}.                         \tag{S2}
\] Indeed reduction modulo \(2\mathbb Z^2\) makes the adjacent squares
central; conversely \([[a,b],c]\) and \([[b,c],a]\) give the two doubled
unit translations, up to signs. Thus \(H_5\) is class two of order 32.
The squared translations are not always independent cyclic subgroups: in
\(H_6\) they coincide and have order two. Formula (S2) uses a lattice
index and does not require independence of those cyclic subgroups.
\(\square\)

\textbf{Lemma UA.SquareRecognition (all finite two-group square markings).}
Every finite two-group string C-group of type $\{4,4\}$ and order $2^n$
is marked-isomorphic to $H_n$; necessarily $n\ge5$.
\begin{proof}
Its two adjacent dihedral parabolics have orders eight and intersect in
order two, so their product already has 32 elements. Thus $n\ge5$.
The complete affine presentation of UA.Affine maps onto the given group $G$;
write $N$ for its kernel. IP gives irredundancy, and UA.Phi identifies
$G/G'$ with the elementary abelian group on its three markers. Summing
those three binary coordinates gives a character sending each reflection
to the nonidentity element of $C_2$. In the affine group this character is
orientation parity. Hence every element of $N$ preserves orientation.

The kernel contains no nonidentity torsion. A finite-order plane isometry
has a fixed point, obtained for example by averaging the orbit of any point.
The square reflection group moves that point into a closed fundamental right
triangle. The stabilizer of a point in that triangle is trivial in the
interior, generated by its wall reflection on an open edge, or the finite
dihedral group generated by the two walls meeting at a vertex. These
stabilizers are conjugates of the proper marked parabolics. Their maps into
$G$ are injective by its exact pair orders and IP, so none contains a
nonidentity element of $N$.

An orientation-preserving plane isometry is a translation or a rotation.
The linear part here has finite order in $D_8$, so a nontranslation rotation
has finite order. Thus all elements of $N$ are translations. Since $N$ is
normal of finite index, it is a $D_8$-invariant sublattice of $\Z^2$, of
index $2^{n-3}$. Under $\Z^2=\Z[i]$, invariance under a quarter-turn makes
it an ideal. The Euclidean algorithm in $\Z[i]$ makes it principal; its
generator has norm a power of two. As $2$ is a unit times $(1+i)^2$,
unique factorization shows that this generator is a unit times a power of
$1+i$. Even powers give the square lattice and odd powers the diagonal
lattice in UA.Square. There is one such ideal at each two-power index.
Therefore $N=\Lambda_n$ and the marked quotient is $H_n$.
\end{proof}

\textbf{Proposition UA.SquareDecomposition (the two square
decompositions).} For every integer \(q\ge1\), let \(M_2(q)\) and
\(M_1(q)\) be the groups obtained by adjoining respectively
\((bcba)^q=1\) and \((cba)^{2q}=1\) to the complete presentation (S1).
They are the square and diagonal lattice quotients, and have the
decompositions (S3) and (S4) below. Here \(D_{2q}\) is the dihedral
group of order \(2q\), including \(D_2=C_2\) and \(D_4=C_2^2\).

\emph{Proof.} Here bcba is translation by \((1,0)\), and \((cba)^2\) is
translation by \((1,1)\). Their signed-permutation orbits therefore
generate exactly the square and diagonal lattices in (S3)--(S4). This
identifies those complete presentations with the indicated lattice
quotients.

In the square case, \(\langle a,c^b\rangle\) acts on the first
coordinate and \(\langle a^b,c\rangle\) on the second. Their
translations have order q; each is \(D_{2q}\), since its reflection has
nonidentity linear part, also when q=1. They commute and intersect
trivially. The reflection b interchanges the two factors and lies
outside their product. They and b generate all three original
generators. Hence \[
 W/q\mathbb Z^2\cong(D_{2q}\times D_{2q})\rtimes C_2.
 \tag{S3}
\]

In the diagonal case, put \(A_0=\langle b^a,b^c\rangle\) and
\(B_0=\langle b,b^{ac}\rangle\). Their reflecting axes are respectively
\(x+y=0,1\) and \(x-y=0,-1\). The first factor has translations along
\((1,1)\), the second along \((1,-1)\), of order q in each case. Each
reflection retains its nonidentity linear part, so the two factors have
order 2q also at q=1. All cross-generators commute. Their linear parts
show that an element of their intersection must be a translation; the
independence of the two diagonal translations modulo the defining
lattice then makes that intersection trivial. Thus
\(A_0B_0\cong D_{2q}\times D_{2q}\).

Both a and c exchange the two factors by conjugation, so normalize their
product. The two individual reflections a,c do not belong to the product
by their linear parts. The only remaining possible nontrivial
intersection with \(\langle a,c\rangle\) is ac. Every translation part
of an element of \(A_0B_0\) with linear part \(-1\) has coordinates of
the same parity, whereas ac has translation part \((0,1)\). The diagonal
lattice preserves this parity condition, so ac is also outside the
product. Finally b already lies in \(B_0\), so the two factors and
\(\langle a,c\rangle\) generate the full group. Consequently \[
 W/q\langle(1,1),(1,-1)\rangle
 \cong(D_{2q}\times D_{2q})\rtimes(C_2\times C_2).
 \tag{S4}
\] These arguments specify the actions by the original affine generators
and prove the decompositions, rather than only equality of orders.
\(\square\)
